\documentclass[a4paper]{article}
\usepackage{amsfonts}
\usepackage{amsmath}
\usepackage{amssymb}
\usepackage{graphicx}     

\begin{document}
  
\noindent \large Auteur: Abdellah El Moussaoui \\
\noindent \large Studierichting: Informatica\\
\noindent \large Oefeningenreeks 6B nr. 17.6\\

\medskip

\normalsize

$s_n =\displaystyle\sum_{k=1}^{n}k^2 \left(k+3\right)$\\

$= \displaystyle\sum_{k=1}^{n} \left(k^3+3k^2\right)$\\

$= \displaystyle\sum_{k=1}^{n}k^3 + 3\displaystyle\sum_{k=1}^{n}k^2$\\

$= \dfrac{n^2\left(n+1\right)^2}{4} + 3\dfrac{n\left(n+1\right)\left(2n+1\right)}{6}$\\

$= \dfrac{n^2\left(n+1\right)^2}{4} + \dfrac{n\left(n+1\right)\left(2n+1\right)}{2}$\\

$= n\left(n+1\right)\left(\dfrac{n\left(n+1\right)}{4} + \dfrac{2n+1}{2}\right)$\\

$= n\left(n+1\right) \left(\dfrac{n\left(n+1\right) + 2\left(2n+1\right)}{4}\right)$\\

$= \dfrac{n\left(n+1\right) \left(n^2+n+4n+2\right)}{4}$\\

$= \dfrac{n\left(n+1\right) \left(n^2+5n+2\right)}{4}$\\

Omdat $n$ en $n+1$ opeenvolgende getallen zijn, is één van beide even.\\

Als $n$ even is, dan is $n^2+5n+2$ ook even.\\

Als $n$ oneven is, dan is $n^2+5n+2$ even.\\

Dus $n^2+5n+2$ is altijd even.\\

$\Rightarrow n\left(n+1\right)\left(n^2+5n+2\right)$ is deelbaar door $4$\\

$\Rightarrow s_n\in\mathbb{N}$\\


\begin{thebibliography}{999}
%\bibitem{a} Referentie
\end{thebibliography}
\end{document} 